% !PRINTED 176
\weberpage{176}
%<<CONT
divisor $S'$ of $S$ of index $j$, but no others, then
$\eta$ satisfies an equation of degree $j$, which is to be regarded as a resolvent
of the form problem. Every function which likewise admits the
permutations of $S'$ is then a rational
function of $\eta$ and of the invariants of the group $S$. The
resolvent of the $\eta$ is a partial or total resolvent, according
as the group $S'$ has a common divisor with the divisors
of $S$ conjugate to it, or is relatively prime
to them (Bd. I, \S.\,156).

\begin{center}
\S.\ 44.

\gesperrt{Klein's Extension of the Fundamental Algebraic Problem.}
\end{center}

The considerations carried out in the preceding section
form a direct generalisation of the \gesperrt{Galois} theory
for a general equation of degree $n$; and this
theory is contained as a special case in the theory of linear substitution groups.

For according to \S.\ 39 the symmetric functions
of $n$ independent variables $x_1, x_2, \ldots, x_n$ are the invariants of the
group $S$ which consists of the permutations of these $n$ variables,
and the equation $\Phi(t) = 0$ is thus for this case, by Bd. I,
\S.\ 145, the \gesperrt{Galois} resolvent of the equation of degree $n$,
whose roots are the quantities $x_i$. We can therefore conceive the general problem of algebra, to solve an equation of degree $n$, as a form problem of a linear substitution group
of dimension $n$. But now there are special equations which can be solved by form problems of lower dimension,
in particular the \gesperrt{pure equations}
which are solvable by a form problem of the first dimension.

For linear homogeneous substitutions of one dimension are
only of the form
\weq{1}{x' = \alpha x,}
and if these are to form a finite group of degree $\mu$,
then the coefficients $\alpha$ must be roots of unity of degree $\mu$.
Conversely, if we let $\alpha$ in (1) run through all the $\mu$th roots of unity,
we have a group of degree $\mu$. This
group has an absolute invariant $x^{\mu}$, and when this is given,
one obtains $x$ as the $\mu$th root of it.
